%! Change in Tandem — Unit 1, Lesson 1.2 — Group Activity (Answer Key)
\documentclass[10pt]{article}
\usepackage{precalculus-article}
\usepackage{precalculus-key}   % pulls in -boxes; do NOT also load -boxes
\newcommand{\TallMath}[1]{$\displaystyle #1\rule[-1.4em]{0pt}{3.2em}$}

\begin{document}

\pageheader{Unit 1, Lesson 1.2}{Group Activity --- Answer Key}

\vspace{0.12in}

\begin{scenariobox}[The Hiker's Day]{plumlight}
\small
A hiker starts up a canyon trail at sunrise. The function $E$ gives her
elevation in meters $t$ hours after starting. She climbs to a peak, descends
to the river at the canyon floor (elevation 0), then climbs the far side.
Every tier reads this same graph --- no drawing required.

\begin{center}
\begin{tikzpicture}[x=0.6cm, y=0.55cm]
  \draw[linegray, very thin, step=1] (0,0) grid (8,5);
  \draw[->, thick] (0,0) -- (8.6,0) node[below right] {\small $t$ (hr)};
  \draw[->, thick] (0,0) -- (0,5.6) node[above] {\small $E(t)$ (m)};
  \foreach \x in {1,2,3,4,5,6,7,8} \node[below] at (\x,0) {\small \x};
  \foreach \y/\lab in {1/100,2/200,3/300,4/400,5/500} \node[left] at (0,\y) {\small \lab};
  \draw[very thick, plum] (0,1) .. controls (1,4) and (2,5) .. (3,5);
  \draw[very thick, plum] (3,5) -- (5,0);
  \draw[very thick, plum] (5,0) .. controls (6.6,0.3) .. (8,3);
  \fill[plum] (0,1) circle (2.5pt);
  \fill[plum] (8,3) circle (2.5pt);
\end{tikzpicture}
\end{center}

The first hours of the climb, in a table:

\begin{center}
\renewcommand{\arraystretch}{1.3}
\begin{tabular}{|c|c|c|c|c|}
\hline
$t$ (hr) & 0 & 1 & 2 & 3 \\
\hline
$E(t)$ (m) & 100 & 300 & 430 & 500 \\
\hline
\end{tabular}
\end{center}
\end{scenariobox}

\vspace{0.1in}

\boxguard
\begin{tcolorbox}[colback=white, colframe=black!40,
    title=\textbf{Tier R --- Remediate},
    fonttitle=\bfseries, arc=2mm, left=3mm, right=3mm, top=2mm, bottom=2mm]

\begin{enumerate}[itemsep=8pt]
  \item Read the graph.

        \begin{enumerate}[label=(\alph*), itemsep=4pt]
          \item $E(0) = $ \ans{100} meters \quad (where the trail starts)
          \item The hiker's highest elevation is \ans{500} meters, reached
                at $t = $ \ans{3} hours.
        \end{enumerate}

  \item Circle one in each pair.

        \begin{enumerate}[label=(\alph*), itemsep=4pt]
          \item For $0 < t < 3$, $E$ is \ans{increasing} /
                \textbf{decreasing}.
          \item For $3 < t < 5$, $E$ is \textbf{increasing} /
                \ans{decreasing}.
        \end{enumerate}

  \item The statement $E(5) = 0$ describes an event on the hike. Write it as
        a full sentence.

        \vspace{0.02in}
        \ansline{Five hours after starting, the hiker's elevation is 0 meters --- she reaches the river at the canyon floor.}

  \item For this function: \quad Input quantity: \ans{time (hours)} \quad
        Output quantity: \ans{elevation (meters)}
\end{enumerate}
\end{tcolorbox}

\vspace{0.1in}

\boxguard
\begin{tcolorbox}[colback=white, colframe=black!40,
    title=\textbf{Tier A --- Approaching Proficiency},
    fonttitle=\bfseries, arc=2mm, left=3mm, right=3mm, top=2mm, bottom=2mm]

\begin{enumerate}[itemsep=8pt]
  \item Name every interval, using inequalities on $t$.

        \vspace{0.05in}
        Increasing for \ans{0} $< t <$ \ans{3} \ and for
        \ans{5} $< t <$ \ans{8} \qquad
        Decreasing for \ans{3} $< t <$ \ans{5}

  \item Use the table of the climb. The hourly gains are
        \ans{200} , \ans{130} , \ans{70} meters. Is the rate of
        change increasing or decreasing? Is the climb concave up or concave
        down?

        \vspace{0.02in}
        \ansline{The gains shrink, so the rate of change is decreasing --- the climb is concave down.}

  \item State the domain and range of $E$ as inequalities.

        \vspace{0.05in}
        Domain: \ans{0} $\le t \le$ \ans{8} \quad
        Range: \ans{0} $\le E \le$ \ans{500}

  \item What is the zero of $E$, and why does a zero here \textit{not} mean
        the hike is over?

        \vspace{0.02in}
        \ansline{The zero is $t = 5$, where $E(5) = 0$. Elevation 0 just means the hiker is at the river's level --- the hike continues up the far side, so a zero output says nothing about the input running out.}
\end{enumerate}
\end{tcolorbox}

\vspace{0.1in}

\boxguard
\begin{tcolorbox}[colback=white, colframe=black!40,
    title=\textbf{Tier E --- Extension},
    fonttitle=\bfseries, arc=2mm, left=3mm, right=3mm, top=2mm, bottom=2mm]

\begin{enumerate}[itemsep=8pt]
  \item On $5 < t < 8$ the hiker climbs slowly at first, then faster and
        faster. Is this stretch concave up or concave down --- and what does
        that say about her pace?

        \vspace{0.02in}
        \ansline{Concave up: the hourly gains are growing, so the rate of change is increasing. Her climbing pace is picking up as she goes.}

  \item Flip the machine: is \textbf{time} a function of \textbf{elevation}
        for this hike? Use the two moments when the hiker is 300 meters up in
        your explanation.

        \vspace{0.02in}
        \ansline{No. The single input ``elevation 300 m'' has more than one output (she is at 300 m at $t = 1$ on the way up and again on the way down) --- one input with two outputs breaks the function rule.}

  \item A second hiker climbs all day, gaining \textit{less and less}
        elevation each hour. Which pair describes his elevation function?

        \begin{enumerate}[label=\Alph*., itemsep=3pt]
          \item Increasing and concave up
          \item Increasing and concave down \quad \textcolor{keyred}{\textbf{$\leftarrow$ correct}}
          \item Decreasing and concave up
          \item Decreasing and concave down
        \end{enumerate}
\end{enumerate}
\end{tcolorbox}

\end{document}
